\documentclass[11pt,a4paper]{article}
\usepackage[utf8]{inputenc}
\usepackage[T1]{fontenc}
\usepackage[spanish]{babel}
\usepackage{geometry,booktabs,hyperref,listings}
\geometry{margin=2.5cm}
\lstdefinelanguage{MiniZinc}{morekeywords={include,int,set,of,array,var,constraint,forall,if,then,else,endif,solve,minimize,output,bool,true,false},sensitive=true,morecomment=[l]{\%},morestring=[b]"}
\lstset{language=MiniZinc,basicstyle=\ttfamily\small,breaklines=true,frame=single,columns=fullflexible}
\title{Optimizacion de modelos MiniZinc para epsilon-consenso}
\author{Proyecto modelo\_simple}
\date{11 de junio de 2026}
\begin{document}
\maketitle

\section{Objetivo}
Se crearon dos variantes nuevas sin sobrescribir los modelos originales:
\begin{itemize}
  \item \texttt{models/simple\_content\_assignment\_epsilon\_consensus\_free\_optimized.mzn}
  \item \texttt{models/simple\_content\_assignment\_epsilon\_consensus\_directed\_optimized.mzn}
\end{itemize}

Las variantes mantienen la semantica central: misma dinamica de opinion, misma restriccion \texttt{alldifferent\_except\_0}, misma monotonocidad de \texttt{max\_diff}, misma condicion final de consenso y mismo objetivo de optimizacion en cada caso.

\section{Contexto: que tipo de problema es}
Estos modelos pertenecen a la familia de problemas de \textit{constraint programming} sobre dominios finitos. En MiniZinc, las decisiones son variables enteras finitas: para cada agente y turno se elige un contenido o cero. Las restricciones describen que asignaciones son validas y la funcion objetivo decide cual solucion es mejor.

En terminos generales, un problema de satisfaccion de restricciones finito puede ser NP-completo. La literatura de CSP muestra que decidir si existe una asignacion que satisfaga restricciones finitas es NP-completo en general, aunque existen subclases polinomiales cuando se restringe la estructura de las relaciones o el grafo de restricciones \cite{csp-complexity,zhuk2017}. MiniZinc tambien se presenta como un lenguaje para problemas combinatorios de alta complejidad y permite resolverlos con CP, MIP, SAT, SMT y otros backends \cite{minizinc-about}.

No es correcto afirmar automaticamente que \textit{este} modelo especifico ya esta formalmente probado NP-completo sin construir una reduccion matematica. Lo honesto es decir:
\begin{itemize}
  \item el modelo es un problema de optimizacion combinatoria sobre dominios finitos;
  \item su espacio bruto de busqueda crece como \((C+1)^{A\cdot Iend}\), antes de propagacion;
  \item combina asignacion multidimensional agente--tiempo--contenido, restricciones globales, dinamica acumulativa y optimizacion;
  \item por contener una estructura de asignacion/scheduling con multiples dimensiones, es razonable esperar comportamiento NP-hard en la practica, como ocurre con muchas variantes de asignacion multidimensional \cite{multidim-assignment}.
\end{itemize}

Por eso, las mejoras no buscan ``volver polinomial'' el problema. Buscan reducir el arbol que el solver necesita explorar: menos simetrias, dominios mas estrechos, busqueda mas informativa y menos salida innecesaria.

\section{Buenas practicas MiniZinc usadas}
El manual de MiniZinc recomienda cotar variables, evitar generadores innecesariamente grandes, usar restricciones globales, agregar restricciones redundantes o de ruptura de simetria cuando aportan propagacion, y experimentar con estrategias de busqueda porque el mejor modelo depende del backend \cite{minizinc-efficient,minizinc-search,minizinc-globals}.

\section{Optimizaciones implementadas}
\subsection{Cotas estrechas}
Originalmente estas cotas estaban fijadas dentro del modelo. En la versi\'on actual se leen desde el archivo \texttt{.dzn}, de modo que la intensidad de intervenci\'on pueda variar entre bater\'ias de prueba sin modificar la formulaci\'on MiniZinc:
\begin{lstlisting}
int: min_zero_per_agent;
int: min_assign_per_agent;
int: max_assign_per_agent = min(C, Iend - min_zero_per_agent);
\end{lstlisting}
Esto reduce dominios: ningun agente puede recibir mas contenidos que \texttt{C}, ni mas que los turnos disponibles despues de exigir ceros. El minimo de asignaciones reales tambien queda parametrizado mediante \texttt{min\_assign\_per\_agent}. En CP, dominios mas pequenos permiten mas poda y, al quedar en los datos, estas cotas se pueden reportar junto con cada benchmark.

\subsection{Ruptura de simetria temporal}
Se agrego:
\begin{lstlisting}
constraint symmetry_breaking_constraint(
  forall(a in AGENT, t in 1..Iend-1)(
    assignment[a,t] = 0 -> assignment[a,t+1] = 0
  )
);
\end{lstlisting}
Esto fuerza que, por agente, los contenidos aparezcan primero y los ceros al final. Es seguro para este modelo porque no hay dinamica externa entre turnos: esperar antes de mostrar el mismo contenido no mejora el estado final. El beneficio es eliminar muchas soluciones equivalentes con ceros intermedios.

\subsection{Busqueda por turnos}
La lista de busqueda ahora es:
\begin{lstlisting}
[assignment[a,t] | t in TIME, a in AGENT]
\end{lstlisting}
Asi el solver decide primero el turno 1 de toda la poblacion, luego el turno 2, etc. Esto alinea la busqueda con \texttt{max\_diff[t]}, que tambien se calcula por turno, y ayuda a detectar antes ramas que no reducen polarizacion.

\subsection{first\_fail e indomain\_median}
Se uso \texttt{first\_fail} para elegir variables con menos valores disponibles e \texttt{indomain\_median} para probar contenidos intermedios primero. En problemas de consenso, los contenidos intermedios suelen actuar como puentes entre extremos.

\subsection{Salida compacta}
Se agrego \texttt{bool: full\_output = false;}. Esto evita imprimir toda la matriz \(A\times Iend\). No cambia la busqueda, pero protege memoria y tiempo de salida en instancias de 100, 200 o 500 agentes.

\section{Cambios descartados}
No se agregaron restricciones por pares de agentes \(O(A^2)\), porque para 500 agentes serian demasiado pesadas. Tampoco se agrego monotonicidad hacia \texttt{target} en el dirigido, porque podria cambiar la semantica original. Se evitaron booleanos auxiliares por celda para no aumentar memoria.

\section{Resultados obtenidos}
Se agregaron dos datos limpios de benchmark:
\begin{itemize}
  \item \texttt{data/benchmark\_epsilon\_free\_5a\_8c\_10i.dzn}
  \item \texttt{data/benchmark\_epsilon\_directed\_5a\_8c\_10i\_target50.dzn}
\end{itemize}

\begin{table}[h]
\centering
\begin{tabular}{llrrrr}
\toprule
Modelo & Solver & Objetivo & Nodos & Fallos & Solve time \\
\midrule
Free original & Chuffed & 10 & 7918 & 4959 & 0.094 s \\
Free optimizado & Chuffed & 10 & 432 & 111 & 0.009 s \\
Directed original & Chuffed & 13 & 882 & 282 & 0.014 s \\
Directed optimizado & Chuffed & 13 & 1749 & 313 & 0.027 s \\
\bottomrule
\end{tabular}
\caption{Benchmarks pequenos con Chuffed y limite de 10 segundos.}
\end{table}

En free, la ruptura temporal redujo drasticamente el arbol. En directed, el original ya era muy competitivo con Chuffed en la instancia pequena; el optimizado conserva salida compacta y simetria temporal, utiles para instancias grandes aunque no siempre ganen en microbenchmarks.

\begin{table}[h]
\centering
\begin{tabular}{llrrrr}
\toprule
Modelo & Solver & Mejor solucion 10s & Nodos & Fallos & Nota \\
\midrule
Free original & Gecode & 14 & 2037624 & 1018801 & Sin prueba de optimalidad \\
Free optimizado & Gecode & 10 & 2097994 & 1048986 & Mejor solucion encontrada \\
Directed original & Gecode & 21 & 1994348 & 997156 & Sin prueba de optimalidad \\
Directed optimizado & Gecode & 20 & 2313502 & 1156734 & Mejor solucion encontrada \\
\bottomrule
\end{tabular}
\caption{Benchmarks pequenos con Gecode y limite de 10 segundos.}
\end{table}

\subsection{Estadisticas por solver en los modelos optimizados}
Tambien se midieron los modelos optimizados contra tres backends disponibles en la instalacion local de MiniZinc: Chuffed, Gecode y OR Tools CP-SAT. La comparacion es util porque cada solver explora el arbol de una forma distinta. Chuffed usa aprendizaje de conflictos; Gecode usa propagacion CP clasica y suele ser mas conservador en memoria; OR Tools CP-SAT traduce el problema a un enfoque SAT/CP entero con propagacion y busqueda propia.

\begin{table}[h]
\centering
\scriptsize
\begin{tabular}{llrrrl}
\toprule
Modelo optimizado & Solver & Objetivo & Fallos & Propagaciones & Solve time \\
\midrule
Free & Chuffed & 10 & 111 & 76156 & 0.009 s \\
Free & Gecode & 10 & 1048986 & 199847374 & 9.896 s, sin prueba completa \\
Free & OR Tools CP-SAT & 10 & 26 & 64409 & 0.159 s \\
Directed & Chuffed & 13 & 313 & 252705 & 0.027 s \\
Directed & Gecode & 20 & 1156734 & 159901470 & 9.903 s, sin prueba completa \\
Directed & OR Tools CP-SAT & 13 & 42 & 138104 & 0.226 s \\
\bottomrule
\end{tabular}
\caption{Estadisticas de los modelos optimizados con distintos solvers. En Gecode, el limite de 10 segundos termino antes de probar optimalidad; por eso se reporta la mejor solucion encontrada.}
\end{table}

La lectura principal es que Chuffed fue el mejor backend para estos benchmarks: probo optimalidad rapido y con pocos fallos. OR Tools CP-SAT tambien probo optimalidad, aunque con mayor tiempo en estas instancias pequenas. Gecode encontro soluciones, pero no alcanzo a cerrar la prueba de optimalidad en 10 segundos. Esto no significa que Gecode sea malo; significa que, para esta familia de instancias, la combinacion de aprendizaje de conflictos de Chuffed parece aprovechar mejor las restricciones reificadas y la estructura de asignacion. Gecode sigue siendo una opcion razonable cuando se quiere evitar que el aprendizaje de nogoods consuma memoria en instancias grandes.

\section{Recomendacion de ejecucion}
Para instancias pequenas y medianas, usar Chuffed:
\begin{lstlisting}
minizinc --solver Chuffed --time-limit 60000 models/simple_content_assignment_epsilon_consensus_free_optimized.mzn data/benchmark_epsilon_free_5a_8c_10i.dzn
minizinc --solver Chuffed --time-limit 60000 models/simple_content_assignment_epsilon_consensus_directed_optimized.mzn data/benchmark_epsilon_directed_5a_8c_10i_target50.dzn
\end{lstlisting}

Como segunda opcion, probar OR Tools CP-SAT:
\begin{lstlisting}
minizinc --solver "OR Tools CP-SAT" --time-limit 60000 models/simple_content_assignment_epsilon_consensus_free_optimized.mzn data/benchmark_epsilon_free_5a_8c_10i.dzn
minizinc --solver "OR Tools CP-SAT" --time-limit 60000 models/simple_content_assignment_epsilon_consensus_directed_optimized.mzn data/benchmark_epsilon_directed_5a_8c_10i_target50.dzn
\end{lstlisting}

Para 200 o 500 agentes, si la RAM crece demasiado, usar Gecode con \texttt{full\_output=false}. Gecode suele aprender menos nogoods que Chuffed y OR Tools CP-SAT, por lo que puede ser mas conservador en memoria, aunque normalmente tarda mas. Para no sobreexplotar RAM: usar limites de tiempo, dejar salida compacta, relajar epsilon al inicio y luego cerrarlo gradualmente.

\begin{thebibliography}{9}
\bibitem{minizinc-efficient} MiniZinc Handbook. \textit{Effective Modelling Practices in MiniZinc}. \url{https://docs.minizinc.dev/en/stable/efficient.html}
\bibitem{minizinc-search} MiniZinc Handbook. \textit{Search}. \url{https://docs.minizinc.dev/en/stable/mzn_search.html}
\bibitem{minizinc-globals} MiniZinc Handbook. \textit{Global constraints}. \url{https://docs.minizinc.dev/en/stable/lib-globals.html}
\bibitem{minizinc-about} MiniZinc overview. \url{https://www.minizinc.org/}
\bibitem{csp-complexity} Complexity of constraint satisfaction. \url{https://en.wikipedia.org/wiki/Complexity_of_constraint_satisfaction}
\bibitem{zhuk2017} Zhuk, D. \textit{A Proof of the CSP Dichotomy Conjecture}. \url{https://arxiv.org/abs/1704.01914}
\bibitem{multidim-assignment} Multidimensional assignment problem. \url{https://en.wikipedia.org/wiki/Multidimensional_assignment_problem}
\end{thebibliography}
\end{document}
